\begin{corollary}[零编码熵并不推出有限层可外置化]\label{cor:conclusion-zero-code-entropy-does-not-imply-finite-layer-externalizability}
设 \(T:X\twoheadrightarrow X\) 为可数集合上的有限纤维满射，并假设对每个 \(k\ge 1\) 都存在长度 \(k\) 的真实向后分支链。若
\[
h_{\mathrm{br}}(T):=
\lim_{\ell\to\infty}\frac{1}{\ell}\log D_\ell(T)=0,
\]
并定义长度-\(\ell\) 单符号精确编码所需的最小字母表大小为
\[
M_\ell(T)
:=
\min\Bigl\{M:\text{存在大小为 }M\text{ 的长度-}\ell\text{ 单符号精确编码}\Bigr\},
\]
则
\[
h_{\mathrm{code}}(T):=
\lim_{\ell\to\infty}\frac{1}{\ell}\log M_\ell(T)=0.
\]
然而，仍不存在任何有限层 exact externalization 可以精确编码全部无限前历史；等价地，精确无限前史外置依然必须使用 \(\omega\)-层 prime-slice 结构。
\end{corollary}
\begin{proof}
由定理 \ref{thm:conclusion-natural-extension-lstep-minimal-alphabet-submultiplicative}，对每个 \(\ell\ge 1\) 都有
\[
M_\ell(T)=D_\ell(T).
\]
因此
\[
\frac{1}{\ell}\log M_\ell(T)=\frac{1}{\ell}\log D_\ell(T),
\]
从而
\[
h_{\mathrm{code}}(T)=h_{\mathrm{br}}(T)=0.
\]

另一方面，由定理 \ref{thm:conclusion-natural-extension-omega-layer-minimal-externalization}，只要系统在每一深度都持续出现真实分支，就不存在任何有限层外置能够精确编码全部无限前历史。故“零编码熵”与“有限层可外置化”在此严格分离。证毕。
\end{proof}

\leanverified{paper\_conclusion\_smith\_p\_kernel\_spectrum\_cone\_simplicial}
\begin{theorem}[Smith \texorpdfstring{$p$}{p}-核谱函数锥的单纯形结构]\label{thm:conclusion-smith-p-kernel-spectrum-cone-simplicial}
固定素数 \(p\)。令 \(\mathcal S_p\) 为一切 Smith \(p\)-核谱函数之集合，即全部可写成
\[
f(k)=\sum_{\ell\ge 1}m_\ell\min(k,\ell),
\qquad
m_\ell\in \ZZ_{\ge 0},
\]
且仅有限多个 \(m_\ell\) 非零的函数 \(f:\ZZ_{\ge 0}\to \ZZ_{\ge 0}\)。再记
\[
g_\ell(k):=\min(k,\ell)
\qquad(\ell\ge 1).
\]
则：
\begin{enumerate}
  \item \(\mathcal S_p\) 在逐点加法下与自由交换幺半群
        \[
        \NN^{(\NN_{\ge 1})}
        \]
        严格同构；
  \item 其有理凸锥
        \[
        \QQ_{\ge 0}\mathcal S_p
        \]
        是单纯形锥；
  \item 其全部极射线恰为
        \[
        \QQ_{\ge 0}\,g_\ell
        \qquad(\ell\ge 1).
        \]
\end{enumerate}
因此，局部 Smith 缺损谱不存在任何非平凡原子关系；任意可加不变量都必在线性上分解到这些高度原子 \(g_\ell\) 上。
\end{theorem}
\begin{proof}
定理 \ref{thm:conclusion-smith-p-kernel-spectrum-discrete-concave-characterization} 与推论 \ref{cor:conclusion-smith-p-kernel-spectrum-free-commutative-monoid} 已给出：每个 \(f\in\mathcal S_p\) 都唯一表示为
\[
f=\sum_{\ell\ge 1}m_\ell g_\ell,
\qquad
m_\ell\in \ZZ_{\ge 0},
\]
且仅有限多个 \(m_\ell\) 非零。故
\[
f\longmapsto (m_\ell)_{\ell\ge 1}
\]
给出 \(\mathcal S_p\cong \NN^{(\NN_{\ge 1})}\) 的交换幺半群同构，第一条成立。

于是 \(\QQ_{\ge 0}\mathcal S_p\) 中的每个元素都唯一写成有限和
\[
\sum_{\ell\ge 1}a_\ell g_\ell,
\qquad
a_\ell\in \QQ_{\ge 0}.
\]
这正是单纯形锥的定义性特征，故第二条成立。

最后，若某个 \(g_\ell\) 不是极射线，则可写成
\[
g_\ell=\sum_{j\ge 1}a_j g_j,
\qquad
a_j\in\QQ_{\ge 0},
\]
且至少有某个 \(j\neq \ell\) 的 \(a_j\) 非零。乘以公共分母后，便得到 \(N g_\ell\) 的两种不同非负整数原子分解，这与上面的唯一性矛盾。故 \(\QQ_{\ge 0}g_\ell\) 必为极射线。反过来，由于所有元素都由这些 \(g_\ell\) 的非负有理线性组合生成，锥的全部极射线只能是这些原子射线。证毕。
\end{proof}

\begin{theorem}[不存在与有限深度截断相容的有限秩加法账本化]\label{thm:conclusion-no-truncation-compatible-finite-rank-additive-ledger-for-exact-history}
设 \(T:X\twoheadrightarrow X\) 为可数集合上的持续分支系统。则不存在精确无限前史编码
\[
\Psi:X^{\leftarrow}\longrightarrow \varprojlim_{\ell\ge 1}(\ZZ\oplus L_\ell),
\]
其中每个 \(L_\ell\) 都是有限生成交换群，并且同时满足下列两条：
\begin{enumerate}
  \item 对每个 \(\ell\ge 1\)，投影
        \[
        \Psi_\ell:=\operatorname{pr}_\ell\circ \Psi
        \]
        精确决定长度-\(\ell\) 的前历史；
  \item 每个 \(\Psi_\ell\) 在相应 prime-slice 标签上都与乘法结构相容地线性化到加法账本
        \[
        \ZZ\oplus L_\ell.
        \]
\end{enumerate}
换言之，精确无限前史码不可能同时满足“每一有限深度都可由有限秩加法账本承载”与“跨深度截断相容”这两项要求。
\end{theorem}
\begin{proof}
反设存在这样的 \(\Psi\)。则对每个固定的 \(\ell\ge 1\)，其截断投影
\[
\Psi_\ell:X^{\leftarrow}\to \ZZ\oplus L_\ell
\]
都会给出一个长度-\(\ell\) 的精确前史外置，并且按假设可在有限生成加法账本 \(\ZZ\oplus L_\ell\) 中完成 prime-slice 线性化。

但定理 \ref{thm:conclusion-primeslice-finite-depth-ledger-incompatibility} 已经排除：任何有限深度的精确 prime-slice 外置，一旦要求压缩进有限秩加法账本，都会与 Gödel 型同态化障碍矛盾。于是每个固定的 \(\ell\) 都已不可能存在这样的 \(\Psi_\ell\)，更不可能存在把它们在全部深度上兼容拼接起来的逆极限编码 \(\Psi\)。故矛盾。证毕。
\end{proof}
